\documentclass[10pt,a4paper]{amsart}
\usepackage[margin=19mm]{geometry}
\usepackage{amsmath,amssymb,mathtools}
\usepackage[scr=rsfs,cal=euler]{mathalfa}
\usepackage{xfrac}
\usepackage[unicode,hidelinks,bookmarks=false]{hyperref}
\newcommand{\ie}{i.e.\ }
\newcommand{\cf}{cf.\ }
\newcommand{\resp}{respectively\ }
\newcommand{\Subsection}{\subsection}
\providecommand{\nobreakdash}{\nobreak\hbox{-}}
\setlength{\emergencystretch}{2em}
\multlinegap0pt
\usepackage{bm}
\usepackage{bbm}
\usepackage{dsfont}
\usepackage{xspace}
\usepackage{cancel}
\usepackage{mathtools}
\usepackage{amsthm}
\makeatletter
\@ifundefined{theorem}{\newtheorem{theorem}{Theorem}}{}
\@ifundefined{corollary}{\newtheorem{corollary}[theorem]{Corollary}}{}
\@ifundefined{proposition}{\newtheorem{proposition}[theorem]{Proposition}}{}
\makeatother
\DeclareMathOperator{\SG}{{\mathbb S}}
\DeclareMathOperator{\im}{im}
\newcommand{\Aut}{{\operatorname{Aut}}}
\newcommand{\F}{\mathcal{F}}
\newcommand{\UAut}{{\operatorname{UAut}}}
\newcommand{\ch}{\operatorname{ch}}
\title[The ${\mathbb S}\_n$-equivariant Euler characteristic of the ]{The ${\mathbb S}\_n$-equivariant Euler characteristic of the moduli space of graphs}
\author{Michael Borinsky, Jos Vermaseren}
\date{Source: arXiv:2306.15598v4 (CC BY 4.0)}
\begin{document}
\maketitle

\noindent\emph{Source and licence.} The ${\mathbb S}_n$-equivariant Euler characteristic of the moduli space of graphs. Version arXiv:2306.15598v4 (CC BY 4.0), distributed under the Creative Commons Attribution 4.0 International licence, as recorded in the arXiv licence metadata for this exact version and shown on the abstract page https://arxiv.org/abs/2306.15598v4. Full rights notice: licenses/arxiv-2306.15598-CC-BY-4.0.txt.\par\smallskip

\noindent\textit{Editorial scope.} Counted unit: the Theorem whose source label is \texttt{thm:equiv} in the article (source0.tex lines 880--890, source label \texttt{thm:equiv}), delivered together with the article's own complete original proof of it (source0.tex lines 891--910, one \texttt{proof} environment of 20 lines). No mathematics is written, repaired or re-derived: statement and proof are the article's own text line for line. Required context retained uncounted: prop:equivC (lines 687--696); cor:chain\_trace (lines 864--869). Every definition, display and label of those lines is retained unchanged. Omitted from this package and not redistributed: 1--686, 697--863, 870--879, 911--1779. Nothing inside a retained excerpt is omitted or altered: every retained body is delivered exactly as the article's lines are written, so no omission receipt of any kind accompanies this package. Citations: the retained excerpts carry no citation command at all, so no bibliography entry of the article is transcribed and no reference is redistributed. Reference adaptation: none was needed and none was made; every reference inside the delivered unit resolves inside it, so no unresolved reference or citation is produced. Printed slips kept literal: no printed word, symbol, hypothesis or number of the counted statement or of its proof was altered. Layout and class: the article's own class is \texttt{amsart} with packages including geometry, fontenc, inputenc, lmodern, pdflscape, diagbox, xcolor, tikz, amsmath, amssymb, amsthm, subcaption, hyperref; the wrapper is the lane's pinned \texttt{amsart} editorial wrapper at 10pt on A4, which restates the theorem-like environments the retained text needs and adds the article's own macro definitions that the retained text needs (\texttt{\textbackslash Aut}, \texttt{\textbackslash F}, \texttt{\textbackslash SG}, \texttt{\textbackslash UAut}, \texttt{\textbackslash ch}, \texttt{\textbackslash im}), with extra wrapper packages (bm, bbm, dsfont, xspace, cancel, mathtools). The delivered numbering is the unit's own. Assets: no retained span contains a figure environment, a graphics-inclusion command, a PDF-page inclusion, a table or a TikZ picture; no image file is copied into this package, the package declares no asset dependency and no image file is redistributed with it. Licence: version arXiv:2306.15598v4 of this article is distributed under the Creative Commons Attribution 4.0 International licence, as recorded in the arXiv licence metadata for this exact version and shown on the abstract page https://arxiv.org/abs/2306.15598v4. A full rights notice is delivered at licenses/arxiv-2306.15598-CC-BY-4.0.txt. Characters: every retained excerpt is pure ASCII, so the delivered engine, XeLaTeX with its default Latin Modern text font, reports no missing character.
\begin{proposition}
\label{prop:equivC}
The $\SG_n$-equivariant Euler characteristic of $\F_{g,n}^\pm$ is given by 
\begin{align*} e_{\SG_n}(\F_{g,n}^\pm) &= \sum_k (-1)^k \ch(\chi_{C_k(\F^\pm_{g,n})}), \end{align*}
where $\ch(\chi_{C_k(\F^\pm_{g,n})})$
is the Fr\"obenius characteristic of $C_k(\F^\pm_{g,n})$ as an $\SG_n$-representation,
\begin{align*} \ch(\chi_{C_k(\F^\pm_{g,n})}) = \frac{1}{n!} \sum_{\pi \in \SG_n} \chi_{C_k(\F^\pm_{g,n})}(\pi) p^\pi, \end{align*}
and $\chi_{C_k(\F^\pm_{g,n})}$ 
is the \emph{character} associated to $C_k(\F^\pm_{g,n})$.
\end{proposition}

\begin{corollary}
\label{cor:chain_trace}
For $g,n \geq 0$ with $2g - 2 + n \geq 0$ and $k \geq 0$, we have %
\begin{align*} \chi_{C_k(\F^{\pm}_{g,n})}(\pi) = \sum_{[G,\Phi]} \frac{1}{|\Aut(G,\Phi)|} \sum_{\alpha \in \pi_G^{-1}(\pi)} \xi^\pm(G,\Phi,\alpha), \end{align*}
where we sum over all isomorphism classes of (not necessarily orientable) connected forested graphs $[G,\Phi]$ of rank $g$, $n$ legs and $k=|\Phi|$ forest edges.
\end{corollary}

\begin{theorem}
For $g,n \geq 0$ with $2g - 2 + n \geq 0$, %
%
%
%
%
\label{thm:equiv}
\begin{align*} e_{\SG_n}(\F^\pm_{g,n}) &= \sum_{[G,\Phi]_U} \frac{(-1)^{|\Phi|}}{|\UAut(G,\Phi)|} \sum_{\alpha \in \UAut(G,\Phi)} \xi^\pm(G,\Phi,\alpha) p^{\pi_G(\alpha)}, \end{align*}
where we sum over all isomorphism classes of connected forested graphs $[G,\Phi]_U$ of rank $g$ with $n$ \emph{unlabeled} legs.

\end{theorem}

\begin{proof}
We can plug the statement of Corollary~\ref{cor:chain_trace} into the definition of the Fr\"obenius characteristic in Proposition~\ref{prop:equivC} to get
\begin{align*} \ch(\chi_{C_k(\F^\pm_{g,n})}) = \frac{1}{n!} \sum_{\pi \in \SG_n} p^\pi \sum_{[G,\Phi]} \frac{1}{|\Aut(G,\Phi)|} \sum_{\alpha \in \pi_G^{-1}(\pi)} \xi^\pm(G,\Phi,\alpha). \end{align*}
%
Next, we can merge the first and the third sum into a summation over the whole group $\UAut(G,\Phi)$, 
as the preimages of different $\pi \in \SG_n$ partition $\UAut(G,\Phi)$:
\begin{align*} \ch(\chi_{C_k(\F^\pm_{g,n})}) = \frac{1}{n!} \sum_{[G,\Phi]} \frac{1}{|\Aut(G,\Phi)|} \sum_{\alpha \in \UAut(G,\Phi)} \xi^\pm(G,\Phi,\alpha) p^{\pi_G(\alpha)} . \end{align*}
Let $L_S(G,\Phi)$ be the set of isomorphism classes of forested graphs $(G,\Phi)$ that only differ by a relabeling of the legs.
%
%
%
%
%
%
Obviously, $\SG_n$ acts transitively on $L_S(G,\Phi)$ by permuting the leg-labels.  The stabilizer of this action is the image of $\pi_G:\UAut(G,\Phi)\rightarrow \SG_n$.
By the orbit-stabilizer theorem, $n! = |\SG_n| = |\im \pi_G| |L_S(G,\Phi)|$.
By the short exact sequence, $1 \rightarrow \Aut(G,\Phi) \rightarrow \UAut(G,\Phi ) \rightarrow \im \pi_G \rightarrow 1$, we have $|\im \pi_G| = |\UAut(G,\Phi)|/ |\Aut(G,\Phi)|$. Hence,
\begin{align*} \ch(\chi_{C_k(\F^\pm_{g,n})}) = \sum_{[G,\Phi]} \frac{1}{|L_S(G,\Phi)| |\UAut(G,\Phi)|} \sum_{\alpha \in \UAut(G,\Phi)} \xi^\pm(G,\Phi,\alpha) p^{\pi_G(\alpha)} . \end{align*}
The terms in this sum do not depend on the leg-labeling of the graphs, so we can just sum over non-leg-labeled graphs and remove the $|L_S(G,\Phi)|$ in the denominator.
\end{proof}

\end{document}
